\chapter{Opportunity Identification in Hydroponic Agriculture}

\section{Purpose of This Session}

In this session, your group will begin identifying a product opportunity within the theme of \textbf{Hydroponics Agriculture}.

\textbf{By the end of this session, your group should:}
\begin{enumerate}
    \item select one product related to hydroponic agriculture;
    \item decide which company your group will represent in developing that product; and
    \item prepare brief notes explaining your choice.
\end{enumerate}

At this stage, the goal is to generate and organize ideas, not to produce a fully validated final answer. Your discussion notes may be written in bullet-point form.

\section{Theme: Hydroponics Agriculture}

All project groups will work within a common theme: \textbf{Hydroponics Agriculture}.

Hydroponics is a method of growing plants without soil, using water mixed with nutrients to support plant growth. A hydroponic system involves not only plant cultivation, but also supporting activities such as nutrient delivery, monitoring, maintenance, harvesting, and handling.

For this project, your group must choose a product that has a \textbf{clear and specific application} in hydroponic agriculture.

\subsection{What Type of Product Can You Choose?}

Your group may choose to design either:
\begin{itemize}
    \item a product \textbf{directly used} in hydroponic agriculture, or
    \item a product that \textbf{indirectly supports} hydroponic agriculture.
\end{itemize}

\textbf{Directly used products} are products involved in the actual operation of a hydroponic system, such as growing, nutrient delivery, or plant support.

\textbf{Indirectly supporting products} are products that assist hydroponic activities in a closely related way, such as maintenance, handling, or system support.

However, the connection to hydroponics must be clear and specific. A product that is too general, or only loosely related to hydroponics, is not suitable.

\subsection{Examples of Possible Product Areas}

The examples below are intended to guide your thinking. They are not the only possible choices.

Your product may relate to one or more of the following areas:
\begin{itemize}
    \item \textbf{Seedling and germination stage} -- for example, seed trays, seed holders, germination pods, or nursery support tools.
    \item \textbf{Nutrient preparation and delivery} -- for example, nutrient mixing containers, dosing devices, pumps, tubing systems, or flow regulators.
    \item \textbf{Plant growth support} -- for example, plant holders, net pots, structural frames, vertical growing racks, or root support components.
    \item \textbf{Environmental monitoring and control} -- for example, pH monitoring devices, water level indicators, temperature sensors, lighting supports, or ventilation-related accessories.
    \item \textbf{Maintenance and cleaning} -- for example, drainage tools, cleaning devices, modular replacement parts, or system inspection aids.
    \item \textbf{Harvesting and post-harvest handling} -- for example, harvesting trays, cutting aids, storage containers, or packaging supports.
\end{itemize}

These examples are only guides. Your chosen product must still be:
\begin{itemize}
    \item clearly related to hydroponic agriculture,
    \item specific rather than too broad,
    \item feasible to discuss and develop within the course, and
    \item suitable as an engineered, discrete, and physical product.
\end{itemize}

\section{In-Class Activities}

\subsection{Activity 1: Choose a Product and Company Context}

The purpose of this activity is to help your group define a clear starting point for the mock product development process.

In this activity, your group must make two decisions:
\begin{enumerate}
    \item What product will your group develop?
    \item Which company will your group represent?
\end{enumerate}

\subsubsection{Part A: Choose a Product}

During your discussion, select one product that your group intends to develop.

When making your decision, consider the following:
\begin{itemize}
    \item \textbf{Market demand} -- Is there a real need or use for the product?
    \item \textbf{Feasibility} -- Is the product realistic to discuss and develop within the course?
    \item \textbf{Group capability} -- Does your group have enough interest, understanding, and ability to work on it?
\end{itemize}

Your selected product must fit the theme of \textbf{Hydroponics Agriculture}. Within that theme, your group may choose any suitable product. The following guidelines should help:

\begin{enumerate}
    \item It is recommended that your group choose a product that members are familiar with, interested in, and motivated to develop.
    \item The product may be either:
    \begin{itemize}
        \item a new idea, or
        \item an improvement or variation of an existing product.
    \end{itemize}
    \item You only need to conduct a simple search to understand whether similar products already exist. You are \textbf{not} required to prove full novelty at this stage.
    \item In many cases, it is easier to choose a product that already exists in the market and then study or improve it.
    \item Your group may discuss the suitability of the product with the instructor, but your group is responsible for choosing a product that is feasible within the course constraints and project requirements.
    \item Although products can take many forms, this activity focuses on \textbf{engineered, discrete, and physical products}.
\end{enumerate}

\subsubsection{Part B: Choose the Company Your Group Will Represent}

After selecting a product, decide which company will develop it.

For this project, each group is assumed to be working within a company. Throughout the development process, you should imagine yourselves as members of that company.

The company context is important because it helps your group make more realistic decisions about:
\begin{itemize}
    \item target users,
    \item available resources,
    \item product positioning,
    \item manufacturing capability, and
    \item market entry strategy.
\end{itemize}

The chosen company may fall into one of the following two scenarios:
\begin{itemize}
    \item a successful, well-established company that already develops similar products; or
    \item a small or new company that plans to enter the market with the product.
\end{itemize}

\textbf{Example:}  
If your group chooses to develop a modular hydroponic growing rack, you may imagine yourselves as:
\begin{itemize}
    \item members of an established agricultural equipment company expanding its hydroponics product line; or
    \item members of a new start-up company introducing an affordable hydroponic system for urban users.
\end{itemize}

\subsubsection{Expected Output for Activity 1}

By the end of Activity 1, each group should prepare a short point-form summary that includes:
\begin{itemize}
    \item the chosen product,
    \item a brief explanation of how it relates to hydroponic agriculture,
    \item the reason your group selected it,
    \item the company your group will represent, and
    \item a brief explanation of why that company is suitable.
\end{itemize}

\subsection{Activity 2: Additional Work for Groups That Want to Go Further}

The following tasks are \textbf{optional for this class session}, but they are \textbf{strongly recommended} because they can help refine your product selection and may serve as a useful foundation for your coursework project.

These tasks do not need to be completed during class, and some may take several days to finish.

\begin{enumerate}
    \item \textbf{Identify customer needs}
    \begin{itemize}
        \item Identify the specific customer needs that your product is intended to address.
        \item Explain how your product will respond to those needs.
    \end{itemize}

    \item \textbf{Outline the product design}
    \begin{itemize}
        \item Identify the main design requirements.
        \item State the key design parameters or performance considerations.
    \end{itemize}

    \item \textbf{Estimate the product structure}
    \begin{itemize}
        \item Estimate the number of parts or components required.
        \item Prepare a rough sketch of the product, if possible.
    \end{itemize}

    \item \textbf{Estimate the expected selling price}
    \begin{itemize}
        \item Consider manufacturing cost,
        \item desired profit margin, and
        \item market competition.
    \end{itemize}
\end{enumerate}

\section{Final Reminder}

At this stage, your group is not expected to have perfect answers. The main goal is to choose a \textbf{clear, suitable, and feasible product opportunity} within hydroponic agriculture, and to place that opportunity within a realistic company context.

A good choice at this stage will make the later stages of the project much easier.